% Interpretation note: modality-specific vs signed cross-modal interactions
\begin{table}[t]
\centering
\small
\caption{How to read cross-modal boards: modality-specific mass vs signed interactions.}
\label{tab:cross-modal-interpretation}
\begin{tabular}{p{2.4cm}p{5.4cm}p{5.0cm}}
\toprule
Layer & What moves & How to read sign / mass \\
\midrule
Modality-specific
  & Additive shift in one block
    (e.g.\ video mean shift)
  & ``Proportion'' $=$ share of detectable
    drift carried by that block alone;
    sign of the mean shift is the direction
    of the modality effect on $X$ (or on $Y$
    if the outcome channel moves). \\
\addlinespace
Pairwise interaction
  & Cross term $a\!\times\!b$ in $X$ and/or $Y$
    (text$\times$audio, text$\times$video, \ldots)
  & \textbf{Positive:} co-movement of the two
    modalities \emph{amplifies} domain gap /
    outcome risk.
    \textbf{Negative:} co-movement \emph{cancels}
    or reverses the gap (antagonistic). \\
\addlinespace
Combo power
  (this board)
  & Same concept drift, different modality
    subsets (T$\times$A vs T$\times$V)
  & Not a signed attribution; compares
    \emph{detectability} when the design
    keeps only those blocks.
    Empirically T$\times$A $>$ T$\times$V for CF/DR. \\
\bottomrule
\end{tabular}
\end{table}
